% Source: https://artofproblemsolving.com/community/c3h1443755_mathcounts_2011_national_sprint_22_27_29_30_geometry
% Stable identity: aops-topic-1443755-q27
Since the circles are externally tangent, $OP=8+2=10$. The radii $OT$ and $PS$ are both perpendicular to the tangent line $TS$, so $OTSP$ is a right trapezoid. Drop a perpendicular from $P$ to $OT$: it has length $TS$, and it cuts off a right triangle with hypotenuse $OP=10$ and leg $8-2=6$. Hence
\[
TS=\sqrt{10^2-6^2}=8 .
\]

\begin{asy}
usepackage("amsmath");
size(6cm);
defaultpen(fontsize(10pt));
path ra(pair P,pair u,pair v,real s=0.3){return P+s*unit(u)--P+s*unit(u)+s*unit(v)--P+s*unit(v);}
pair O=(0,0),P=(10,0),T=(4.8,6.4),Sx=(11.2,1.6),Q=(8,0);
draw(circle(O,8)); draw(circle(P,2));
draw((T-0.25*(Sx-T))--(Sx+0.35*(Sx-T)),linewidth(0.9));
draw(O--T,gray); draw(P--Sx,gray); draw(O--P,gray);
draw(O--Sx,red+dashed);
draw(ra(T,O-T,Sx-T,0.6)); draw(ra(Sx,P-Sx,T-Sx,0.35));
dot(O);dot(P);dot(T);dot(Sx);dot(Q);
label("$O$",O,SW); label("$P$",P,SE); label("$T$",T+0.9*unit(T)); label("$S$",Sx+1.0*(0.6,0.8)); label("$Q$",Q,1.3*SW);
label("$8$",(2.4,3.2),NW); label("$2$",(P+Sx)/2+(0.45,-0.3));
label("\mbox{}",(min(currentpicture,user=true)+max(currentpicture,user=true))/2,black); // keeps the text after the figure black when LaTeX typesets the labels inline
\end{asy}

Now triangle $OTS$ has a right angle at $T$, so
\[
OS=\sqrt{OT^2+TS^2}=\sqrt{64+64}=\boxed{8\sqrt{2}}\ \text{units}.
\]
